\documentclass[11pt,a4paper]{article}
\usepackage{fontspec}
\usepackage{amsmath,amssymb,amsthm}
\usepackage{mathtools}
\usepackage{bm}
\usepackage{booktabs}
\usepackage{geometry}
\usepackage{enumitem}
\usepackage{xcolor}
\usepackage[pdfencoding=auto,psdextra]{hyperref}

\geometry{margin=2.5cm}

\newcommand{\E}{\mathbb{E}}
\newcommand{\R}{\mathbb{R}}
\newcommand{\KL}{\mathrm{KL}}
\newcommand{\N}{\mathcal{N}}
\newcommand{\sg}{\mathrm{sg}}
\newcommand{\Xspace}{\mathcal{X}}
\newcommand{\Za}{\mathcal{Z}_{A}}
\newcommand{\Zb}{\mathcal{Z}_{B}}
\newcommand{\Enc}{\mathcal{E}}
\newcommand{\Dec}{\mathcal{D}}
\newcommand{\Tmap}{\mathsf{T}}
\newcommand{\Tinv}{\mathsf{T}^{-1}}
\newcommand{\push}{\#}
\newcommand{\Jac}{\mathbf{J}}
\newcommand{\norm}[1]{\left\lVert #1 \right\rVert}
\newcommand{\score}{\mathbf{s}}

\theoremstyle{definition}
\newtheorem{definition}{Definition}
\newtheorem{proposition}{Proposition}
\newtheorem{remark}{Remark}
\newtheorem{observation}{Observation}

\title{\textbf{Avoiding the Student$\to$Teacher Inverse in Cross-Latent On-Policy Distillation}\\[0.3em]
\large Working Analysis (Draft)}
\author{Jayce-Ping}
\date{}

\begin{document}
\maketitle

\begin{abstract}
This is a working (temporary) companion note to
\texttt{cross\_latent\_distillation\_problem.tex}. In the cross-latent setting with
$d_A>d_B$ (the teacher VAE has more channels than the student VAE), the forward transport
$\Tmap:\Za\to\Zb$ (teacher$\to$student, high$\to$low) is easy and well-learned, whereas the
inverse $\Tinv:\Zb\to\Za$ (student$\to$teacher, low$\to$high) is underdetermined and hard.
We ask whether on-policy distillation can be carried out \emph{using only the well-learned
forward map $\Tmap$}, bypassing $\Tinv$. The answer is yes, and the reason is structural:
the inverse is required \emph{only} because we insist on a \emph{pointwise vector-field}
target evaluated at student-chosen states. Measures push forward canonically under any
measurable map, and differentiation of losses built from forward maps uses their
\emph{adjoints} (vector--Jacobian products), which always exist. Reformulating the objective
in either of these two ways removes $\Tinv$ entirely. We give the precise argument and three
concrete routes, and map them onto the existing L0/L1 structure of XOPD.
\end{abstract}

\section{Setup and Notation}
\label{sec:setup}

We reuse the notation of the parent note. Two VAEs induce latent spaces
$\Za\subseteq\R^{d_A}$, $\Zb\subseteq\R^{d_B}$ with encoders/decoders
$(\Enc_A,\Dec_A)$, $(\Enc_B,\Dec_B)$ over a shared pixel space $\Xspace$. The teacher is
frozen with velocity $v_A$ on $\Za$; the student has trainable velocity $v_\theta$ on $\Zb$.
Empirically
\[
d_A > d_B ,
\]
so the forward transport
\[
\Tmap:\Za\to\Zb\quad(\text{teacher}\to\text{student, high}\to\text{low})
\]
is a submersion (many-to-one), easy and well-learned; its would-be inverse
\[
\Tinv:\Zb\to\Za\quad(\text{student}\to\text{teacher, low}\to\text{high})
\]
is a dimensional expansion, underdetermined and hard. On-policy distillation evaluates its
loss on the student-visited distribution $\rho_\theta$ (the marginal of a student rollout
$\{z_B(t)\}$ under $v_\theta$).

\paragraph{Where the inverse enters.}
The pushforward velocity field is, formally,
\begin{equation}
(\Tmap\push v_A)(z_B,t)
\;=\;
\Jac_{\Tmap}\!\big(\Tinv(z_B)\big)\,v_A\!\big(\Tinv(z_B),t\big).
\label{eq:pushforward}
\end{equation}
Using \eqref{eq:pushforward} as a \emph{pointwise target} at a student state $z_B\sim\rho_\theta$
requires a preimage $\Tinv(z_B)$: this, and only this, is what forces the hard inverse.

\section{Why the Inverse Is Hard, and When It Is Truly Needed}
\label{sec:why}

\begin{definition}[Fiber and section]
\label{def:section}
For $z_B\in\Zb$ the \emph{fiber} is $\mathcal{F}(z_B)\coloneqq\Tmap^{-1}(z_B)=\{z\in\Za:\Tmap(z)=z_B\}$,
a $(d_A-d_B)$-dimensional set when $\Tmap$ is a full-rank submersion. A \emph{section} of $\Tmap$
is a right inverse $s:\Zb\to\Za$ with $\Tmap\circ s=\mathrm{id}_{\Zb}$; equivalently $s(z_B)\in\mathcal{F}(z_B)$.
Learning $\Tinv$ is exactly learning a section.
\end{definition}

\begin{proposition}[Vector-field pushforward is generally ill-defined under a non-injective map]
\label{prop:projectability}
Let $\Tmap$ be a submersion with $d_A>d_B$. The expression
$z_B\mapsto \Jac_{\Tmap}(z)\,v_A(z)$ is independent of the choice of representative
$z\in\mathcal{F}(z_B)$ \emph{if and only if} $v_A$ is $\Tmap$-projectable, i.e.\ the value
$\Jac_{\Tmap}(z)v_A(z)$ is constant along each fiber. In general it is not; hence
$\Tmap\push v_A$ is \emph{not} a well-defined vector field, and a pointwise target
\eqref{eq:pushforward} is only defined \emph{after} fixing a section $s=\Tinv$.
\end{proposition}

\begin{remark}[The real difficulty is the formulation, not network capacity]
\label{rem:formulation}
Prop.~\ref{prop:projectability} says the pointwise velocity target is ill-posed \emph{as a
mathematical object}: different sections give different targets, and the ``correct'' section
must recover teacher-exclusive fiber coordinates that $z_B$ does not determine
(the underdetermination of the parent note). Enlarging the network for $\Tinv$ therefore has
diminishing returns; the leverage is in changing the objective so that no section is required.
\end{remark}

\begin{proposition}[Measures always push forward]
\label{prop:measure}
For \emph{any} measurable $\Tmap:\Za\to\Zb$, the pushforward measure $\Tmap\push p_A$, defined
by $(\Tmap\push p_A)(E)=p_A(\Tmap^{-1}(E))$, exists and is unique. In particular the
projected-teacher marginal $\Tmap\push p_A^{0}$ (and its noised versions along the flow) is
well-defined with no choice of section.
\end{proposition}

\begin{observation}[The escape principle]
\label{obs:principle}
Combining Prop.~\ref{prop:projectability} and Prop.~\ref{prop:measure}: replace the
\emph{pointwise vector-field} target by a \emph{distributional} target (which pushes forward
canonically under $\Tmap$), and/or measure the discrepancy in a space reachable by
\emph{forward} maps so that gradients use adjoints (VJPs). Both differentiation and expectation
then use only $\Tmap$ (and transposes of forward maps), never the set-inverse $\Tinv$.
\end{observation}

\section{Route A: Distribution / Score Matching (DMD/VSD style)}
\label{sec:routeA}

Replace the pointwise L1 target by matching the student's on-policy marginal to the projected
teacher marginal. Let $q_t\coloneqq(\Tmap\push p_A)(\cdot,t)$ (well-defined by
Prop.~\ref{prop:measure}) and let $\rho_{\theta,t}$ be the student's rollout marginal at time
$t$. Minimize a reverse-KL / score objective
\begin{equation}
\min_\theta\ \E_t\,\big[\,w(t)\,\KL\!\big(\rho_{\theta,t}\,\big\|\,q_t\big)\big],
\qquad
\nabla_\theta \approx
\E_t\,\E_{z_B\sim\rho_{\theta,t}}\!\Big[\big(\score_{\theta}(z_B,t)-\score_{q}(z_B,t)\big)\,\partial_\theta z_B\Big],
\label{eq:dmd}
\end{equation}
where $\score_{\theta}=\nabla\log\rho_{\theta,t}$ (the ``fake'' score) and
$\score_{q}=\nabla\log q_t$ (the ``real''/target score). Both scores are evaluated
\emph{at student states $z_B\in\Zb$}: no inverse. The target score $\score_q$ is obtained by
fitting a small student-space score model on samples $z_B^0=\Tmap(z_A^0)$, $z_A^0\sim p_A^0$
(teacher samples pushed forward through the \emph{easy} $\Tmap$).
\begin{itemize}[leftmargin=1.4em,itemsep=2pt]
  \item \textbf{On-policy:} yes ($z_B\sim\rho_\theta$).
  \item \textbf{Uses $\Tinv$:} no. Only forward $\Tmap$ (to build target samples).
  \item \textbf{Cost:} one auxiliary score network; adversarial-style two-timescale training.
  \item \textbf{Why principled:} sidesteps Prop.~\ref{prop:projectability} entirely by targeting
        a measure rather than a vector field.
\end{itemize}

\section{Route B: Shared Pixel Space via Adjoints (VJP, not inverse)}
\label{sec:routeB}

Both decoders map into the shared pixel space $\Xspace$. Measure the on-policy discrepancy after
a \emph{forward} decode of the student's clean estimate
$\hat z_B^{0}(z_B,t)$, with a stop-gradient teacher target:
\begin{equation}
L(\theta)=\E_{z_B\sim\rho_\theta}\Big\|\,\Dec_B\big(\hat z_B^{0}\big)
-\sg\!\big[\,x^\star(z_B,t)\,\big]\Big\|^2,
\qquad
x^\star \coloneqq \Dec_A\big(\Psi_A^{0}(\Enc_A(\Dec_B(\hat z_B^{0})),t)\big),
\label{eq:pixel}
\end{equation}
where $\Psi_A^{0}$ is the teacher's clean-sample estimator (one teacher query). The teacher
target uses the \emph{frozen} composition $\Enc_A\circ\Dec_B$; this is \emph{not} a learned
low$\to$high network---the encoder $\Enc_A$ legitimately produces $d_A$ channels from the
decoded image. Crucially, the $\theta$-gradient flows back only through the forward map $\Dec_B$
via its \emph{vector--Jacobian product} $\Jac_{\Dec_B}^{\!\top}$, which always exists; nowhere is
$\Tmap$ inverted.
\begin{itemize}[leftmargin=1.4em,itemsep=2pt]
  \item \textbf{On-policy:} yes.
  \item \textbf{Uses $\Tinv$:} no. Only forward/frozen maps and their adjoints.
  \item \textbf{Cost:} a decode+encode through two VAEs per step; round-trip VAE loss $\delta$.
  \item \textbf{Why principled:} an $x_0$ (clean-sample) target avoids the projectability issue of
        Prop.~\ref{prop:projectability}; differentiation uses $\Jac^{\top}$ (adjoint), not the
        set-inverse.
\end{itemize}

\section{Route C: A $\Tmap$-Anchored Minimal-Norm Section}
\label{sec:routeC}

If a pointwise velocity target is desired, do \emph{not} learn a free $\Tinv$. Anchor it to the
already-good $\Tmap$ by the minimal-norm section (Moore--Penrose pseudo-inverse in the affine
case $\Tmap=Az+b$):
\begin{equation}
\Tinv(z_B)\;\coloneqq\;A^{+}(z_B-b),
\qquad A A^{+}=\mathrm{id}_{\Zb},
\label{eq:pinv}
\end{equation}
subject only to the right-inverse constraint $\Tmap\circ\Tinv=\mathrm{id}$ (not a full cycle).
The teacher-exclusive fiber coordinates are set to their conditional mean; a small learned
residual, optionally conditioned on the student DiT hidden states $h_\theta$ (parent note,
\S ``Conditional Pullback''), refines them:
$\Tinv_c(z_B)=A^{+}(z_B-b)+g_\phi(z_B,h_\theta)$ with $g_\phi$ zero-initialized.
\begin{itemize}[leftmargin=1.4em,itemsep=2pt]
  \item \textbf{On-policy:} yes; keeps the existing pointwise formulation.
  \item \textbf{Uses $\Tinv$:} yes, but reduced from ``a free low$\to$high network'' to
        ``$\Tmap$-anchored pseudo-inverse $+$ small residual''.
  \item \textbf{Cost:} minimal; reuses $\Tmap$'s parameters.
\end{itemize}

\section{Mapping onto XOPD's L0/L1}
\label{sec:l0l1}

\begin{itemize}[leftmargin=1.4em,itemsep=3pt]
  \item \textbf{L0 (clean-sample regression) is already inverse-free:} it only needs sample
        transport $z_B^0=\Tmap(z_A^0)$ through the easy forward $\Tmap$.
  \item \textbf{L1 (on-policy transition / mean matching) is where $\Tinv$ bites:} it needs the
        teacher's opinion at student states. Replace it by Route~A (student-space score
        difference) or Route~B (pixel $x_0$ matching) to make L1 depend only on the forward
        $\Tmap$; or use Route~C to keep L1 pointwise but with a cheap anchored inverse.
\end{itemize}

\begin{table}[h]
\centering
\small
\begin{tabular}{lccll}
\toprule
Route & On-policy & Uses $\Tinv$ & Key cost & Principle \\
\midrule
A: score / DMD          & yes & no  & aux.\ score net, 2-timescale & measure pushforward \\
B: pixel adjoint        & yes & no  & 2$\times$VAE per step, $\delta$ & VJP of forward maps \\
C: anchored pseudo-inv. & yes & yes (cheap) & minimal & $\Tmap$-anchored section \\
\bottomrule
\end{tabular}
\caption{Three ways to reduce or remove the hard student$\to$teacher inverse in on-policy L1.}
\end{table}

\section{Recommendation}
\label{sec:reco}

\begin{itemize}[leftmargin=1.4em,itemsep=2pt]
  \item \textbf{Minimal change, runs now:} Route~C (pseudo-inverse $+$ hidden-state residual),
        keeping the current pointwise formulation.
  \item \textbf{Cleanest, truly eliminates the inverse:} Route~A (score/DMD), which converts the
        ill-posed vector-field-inversion into a well-posed measure-matching problem while staying
        on-policy.
  \item \textbf{Middle ground:} Route~B, on-policy and inverse-free, at the price of per-step VAE
        round-trips.
\end{itemize}

\paragraph{One-line takeaway.}
The inverse $\Tinv$ is required only because we chose a \emph{pointwise vector-field} target at
student-chosen states. Switching to distributional / score / $x_0$ objectives---which push
forward canonically under the easy $\Tmap$ and whose gradients use only adjoints (VJPs) of
forward maps---removes $\Tinv$ from the problem.

\end{document}
